
%%%%%%%%%%%%%%%%%%%%%%%%%%%%%%%%%%%%%%%%%%%%%%%%%%%%%%%%%%%%%%%%%%%%%%%%%%%%%%%%
%%%% DOCUMENT SETUP                                                         %%%%
%%%%%%%%%%%%%%%%%%%%%%%%%%%%%%%%%%%%%%%%%%%%%%%%%%%%%%%%%%%%%%%%%%%%%%%%%%%%%%%%
%\usepackage[showframe]{geometry}
\usepackage[pdfusetitle]{hyperref}
\usepackage{placeins}

\usepackage[titles]{tocloft}
\renewcommand\contentsname{}


% Reduce spacing around displaymath; increased line spacing.
\usepackage[nodisplayskipstretch]{setspace}
\setstretch{1.1}
\setlength{\parskip}{0.5em}
\setlength{\parindent}{0pt}


% Good footnotes.
\usepackage[hang,flushmargin,multiple,bottom,perpage]{footmisc}
\setlength\footnotemargin{1em}
\interfootnotelinepenalty=10000


% Reduce spacing before and after sub/section headers.
\usepackage[medium]{titlesec}
\titlespacing{\section}{0pt}{0.25\baselineskip}{0.25\baselineskip}
\titlespacing{\subsection}{0pt}{0.25\baselineskip}{-0.25\baselineskip}
\titlespacing{\subsubsection}{0pt}{0pt}{0pt}

\usepackage{titling, keyval}

% Reduce space before title.
\setlength{\droptitle}{-5em}

% Left-align title, authors, date.
\pretitle{\begin{flushleft}\LARGE}
\posttitle{\par\end{flushleft}\vspace{-1em}}

\preauthor{\begin{flushleft}}
\postauthor{\end{flushleft}}

\predate{\begin{flushleft}\small\it}
\postdate{\end{flushleft}\vspace{-3em}}


% Custom authorship command, since we have many authors.
\newcommand{\correspondence}{$\vcenter{\hbox{\footnotesize\Letter}}$\xspace}
\newcommand{\authorship}[2]{%
	{\setstretch{1} #1 \\[-1pt] {\small\textit{\textsc{#2}}} \\[1.25em]}
}

% No abstract title.
\renewcommand{\abstractname}{}




% Lets us literally flag a paragraph.
\usepackage{tabto, marginfix, fontawesome}
\def\flag{\protect\flag}
\def\flag{\tabto*{-0.25in}\makebox[0cm]{{\color{black}\normalfont\color{red}\faFlag}}\tabto*{\TabPrevPos}}


% Annotations and definitions.
\usepackage{cuted, xparse}
\setlength\marginparwidth{1.5in}
\newcommand{\annotate}[1]{\marginpar{\footnotesize\color{gray}#1}}
\NewDocumentCommand{\define}{m o}{\emph{#1}\annotate{\IfValueTF{#2}{#2}{#1}}}


% Sub-enumerates include the enumerate.
\usepackage{enumitem}
\setlist[enumerate,2]{label*=\arabic*.}



%%%%%%%%%%%%%%%%%%%%%%%%%%%%%%%%%%%%%%%%%%%%%%%%%%%%%%%%%%%%%%%%%%%%%%%%%%%%%%%%
%%%% FONTS                                                                  %%%%
%%%%%%%%%%%%%%%%%%%%%%%%%%%%%%%%%%%%%%%%%%%%%%%%%%%%%%%%%%%%%%%%%%%%%%%%%%%%%%%%
\usepackage[cal=txupr]{mathalpha}
\usepackage{newtxtext, dsfont}
\usepackage[nopatch=footnote]{microtype}
\usepackage[misc]{ifsym}



%%%%%%%%%%%%%%%%%%%%%%%%%%%%%%%%%%%%%%%%%%%%%%%%%%%%%%%%%%%%%%%%%%%%%%%%%%%%%%%%
%%%% VISUALS                                                                %%%%
%%%%%%%%%%%%%%%%%%%%%%%%%%%%%%%%%%%%%%%%%%%%%%%%%%%%%%%%%%%%%%%%%%%%%%%%%%%%%%%%
\usepackage{quiver, caption, tikz, xcolor, wrapfig}
\usetikzlibrary{graphs, calc, decorations, decorations.markings, patterns, patterns.meta}


% Captions.
\captionsetup{format=plain, font={it, footnotesize}}
\setlength{\belowcaptionskip}{-1em}


% IBM Design Library palette. (Colorblind-safe) %% https://tinyurl.com/27azp7vh
\definecolor{ibm-gold}{HTML}{ffb000}
\definecolor{ibm-orange}{HTML}{fe6100}
\definecolor{ibm-pink}{HTML}{dc267f}
\definecolor{ibm-violet}{HTML}{785ef0}
\definecolor{ibm-blue}{HTML}{648fff}


% Paul Tol's high-contrast palette. Colors remain distinct in monochrome. %% https://tinyurl.com/2b9mab2z
\definecolor{tol-hc-yellow}{RGB}{221,170,51}
\definecolor{tol-hc-red}{RGB}{187,85,102}
\definecolor{tol-hc-blue}{RGB}{0,68,136}


% Paul Tol's medium-contrast palette. Colors remain distinct in monochrome. (Colorblind-safe) %% https://tinyurl.com/2b9mab2z
\definecolor{tol-mc-lightyellow}{RGB}{238,204,102}
\definecolor{tol-mc-lightred}{RGB}{238,153,170}
\definecolor{tol-mc-lightblue}{RGB}{102,153,204}
\definecolor{tol-mc-darkyellow}{RGB}{153,119,0}
\definecolor{tol-mc-darkred}{RGB}{153,68,85}
\definecolor{tol-mc-darkblue}{RGB}{0,68,136}


% Paul Tol's vibrant palette. (Colorblind-safe) %% https://tinyurl.com/2b9mab2z
\definecolor{tol-v-blue}{RGB}{0,119,187}
\definecolor{tol-v-cyan}{RGB}{51,187,238}
\definecolor{tol-v-teal}{RGB}{0,153,136}
\definecolor{tol-v-orange}{RGB}{238,119,51}
\definecolor{tol-v-red}{RGB}{204,51,17}
\definecolor{tol-v-magenta}{RGB}{238,51,119}
\definecolor{tol-v-grey}{RGB}{187,187,187}


% Okabe and Ito's palette. (Colorblind-safe) %% https://tinyurl.com/22lwpexl
\definecolor{oi-orange}{RGB}{230,159,0}
\definecolor{oi-cyan}{RGB}{86,180,233}
\definecolor{oi-green}{RGB}{0,158,115}
\definecolor{oi-yellow}{RGB}{240,228,66}
\definecolor{oi-blue}{RGB}{0,114,178}
\definecolor{oi-red}{RGB}{213,94,0}
\definecolor{oi-pink}{RGB}{204,121,167}

% ArXiV color (RIP).
\definecolor{arxiv}{HTML}{AE2A24}


% Define a color theme; use Paul Tol's dark, high-contrast palette.
\def\primarycolor{tol-mc-darkred}
\def\secondarycolor{tol-mc-lightblue}
\def\tertiarycolor{tol-mc-lightred}


%%%%%%%%%%%%%%%%%%%%%%%%%%%%%%%%%%%%%%%%%%%%%%%%%%%%%%%%%%%%%%%%%%%%%%%%%%%%%%%%
%%%% ENVIRONMENTS                                                           %%%%
%%%%%%%%%%%%%%%%%%%%%%%%%%%%%%%%%%%%%%%%%%%%%%%%%%%%%%%%%%%%%%%%%%%%%%%%%%%%%%%%
\usepackage{amsthm, thmtools}

% Define standard heights, widths, etc.
\def\spaceabove{0.5\parskip}		% thmtools above
\def\spacebelow{0.25\parskip}		% thmtools below

\def\innerleftmargin{0.5em}			% mdframed margins
\def\innerrightmargin{0.5em}
\def\innertopmargin{0.5em}
\def\innerbottommargin{0.25em}

\def\postheadhook{\strut}
\def\prefoothook{\looseness=-1\strut}

\def\rulewidth{1pt}					% mdframed rule width


\declaretheoremstyle[
	spaceabove=\spaceabove,
	spacebelow=\spacebelow,
	bodyfont=\itshape
]{compact}


\declaretheoremstyle[
	spaceabove=\spaceabove,
	spacebelow=\spacebelow,
	bodyfont=\normalfont,
	mdframed={
		innerleftmargin=\innerleftmargin,
		innerrightmargin=\innerrightmargin,
		innertopmargin=\innertopmargin,
		innerbottommargin=\innerbottommargin,
		backgroundcolor=gray!15,
		linewidth=\rulewidth
	}
]{shaded}

\declaretheoremstyle[
	spaceabove=\spaceabove,
	spacebelow=\spacebelow,
	bodyfont=\itshape,
	mdframed={
		innerleftmargin=\innerleftmargin,
		innerrightmargin=\innerrightmargin,
		innertopmargin=\innertopmargin,
		innerbottommargin=\innerbottommargin,
		backgroundcolor=gray!15,
		linewidth=\rulewidth
	}
]{shaded-it}


\declaretheoremstyle[
	spaceabove=\spaceabove,
	spacebelow=\spacebelow,
	bodyfont=\normalfont,
	mdframed={
		innerleftmargin=\innerleftmargin,
		innerrightmargin=\innerrightmargin,
		innertopmargin=\innertopmargin,
		innerbottommargin=\innerbottommargin,
		backgroundcolor=gray!15,
		linewidth=\rulewidth,
		nobreak
	}
]{shaded-nobreak}

\declaretheoremstyle[
	spaceabove=\spaceabove,
	spacebelow=\spacebelow,
	mdframed={
		innerleftmargin=\innerleftmargin,
		innerrightmargin=\innerrightmargin,
		innertopmargin=\innertopmargin,
		innerbottommargin=\innerbottommargin,
		backgroundcolor=gray!15,
		linewidth=\rulewidth,
		nobreak
	}
]{shaded-nobreak-it}


\declaretheoremstyle[
	spaceabove=\spaceabove,
	spacebelow=\spacebelow,
	bodyfont=\normalfont,
	mdframed={
		innerleftmargin=\innerleftmargin,
		innerrightmargin=\innerrightmargin,
		innertopmargin=\innertopmargin,
		innerbottommargin=\innerbottommargin,
		backgroundcolor=almond!35,
		linewidth=\rulewidth
	}
]{shaded-almond}


\declaretheoremstyle[
	spaceabove=\spaceabove,
	spacebelow=\spacebelow,
	bodyfont=\normalfont,
	mdframed={
		innerleftmargin=\innerleftmargin,
		innerrightmargin=\innerrightmargin,
		innertopmargin=\innertopmargin,
		innerbottommargin=\innerbottommargin,
		backgroundcolor=almond!35,
		linewidth=\rulewidth,
		nobreak,
	}
]{shaded-almond-nobreak}


\declaretheoremstyle[
	spaceabove=\spaceabove,
	spacebelow=\spacebelow,
	bodyfont=\normalfont,
	mdframed={
		innerleftmargin=\innerleftmargin,
		innerrightmargin=\innerrightmargin,
		innertopmargin=\innertopmargin,
		innerbottommargin=\innerbottommargin,
		backgroundcolor=\secondarycolor!15,
		linewidth=\rulewidth,
		nobreak
	}
]{shaded-blue-nobreak}

\declaretheoremstyle[
	spaceabove=\spaceabove,
	spacebelow=\spacebelow,
	bodyfont=\normalfont,
	mdframed={
		innerleftmargin=\innerleftmargin,
		innerrightmargin=\innerrightmargin,
		innertopmargin=\innertopmargin,
		innerbottommargin=\innerbottommargin,
		backgroundcolor=\primarycolor!15,
		linewidth=\rulewidth,
		nobreak
	}
]{shaded-red-nobreak}

\declaretheoremstyle[
	spaceabove=\spaceabove,
	spacebelow=\spacebelow,
	bodyfont=\normalfont,
	mdframed={
		innerleftmargin=\innerleftmargin,
		innerrightmargin=\innerrightmargin,
		innertopmargin=\innertopmargin,
		innerbottommargin=\innerbottommargin,
		backgroundcolor=tol-v-teal!15,
		linewidth=\rulewidth,
		nobreak
	}
]{shaded-teal-nobreak}


\declaretheorem[%
%	postheadhook=\postheadhook,
%	prefoothook=\prefoothook,
	numberwithin=subsection,
	style=compact
]{theorem}

\declaretheorem[%
%	postheadhook=\postheadhook,
%	prefoothook=\prefoothook,
	sharenumber=theorem,
	style=compact
]{lemma, definition, corollary, proposition}

\declaretheorem[%
	postheadhook=\postheadhook,
	prefoothook=\prefoothook,
	sharenumber=theorem,
	style=shaded-nobreak-it,
	title=Theorem
]{bigtheorem}

\declaretheorem[%
	postheadhook=\postheadhook,
	prefoothook=\prefoothook,
	sharenumber=theorem,
	style=shaded-nobreak-it,
	title=Definition
]{bigdefinition}

\declaretheorem[%
	postheadhook=\postheadhook,
	prefoothook=\prefoothook,
	sharenumber=theorem,
	style=shaded-nobreak
]{example, conjecture}

\declaretheorem[%
	postheadhook=\postheadhook,
	prefoothook=\prefoothook,
%	numberwithin=section,
	style=shaded
]{question}

\declaretheorem[%
	postheadhook=\postheadhook,
	prefoothook=\prefoothook,
%	sharenumber=question,
	style=shaded-almond
]{goal}

\declaretheorem[%
	postheadhook=\postheadhook,
	prefoothook=\prefoothook,
	numbered=no,
	style=shaded-blue-nobreak,
	title={Author's Note}
]{authorsnote}

\declaretheorem[%
	postheadhook=\postheadhook,
	prefoothook=\prefoothook,
	numbered=no,
	style=shaded-red-nobreak,
	title={Reader's Note}
]{readersnote}


% Algorithm environment.
\usepackage[normalem]{ulem}
\usepackage[many]{tcolorbox}
\newtcolorbox[auto counter]{boxedalgorithm}[3]{%
	label=#2,
	%
	enhanced,
	sharp corners,
	frame hidden,
	%
	center,
	width=0.75\linewidth,
	top=0.5em,
	bottom=0.5em,
	left=0.5em,
	right=0.5em,
	%
	before skip=1em,
	after skip=1.5em,
	%
	attach title to upper={#3},
	coltitle=black,
	title={\uline{\textbf{Algorithm \thetcbcounter}{#1}\hfill}},
	%
	breakable
}

\usepackage[many]{tcolorbox}
\newtcolorbox[auto counter]{headlessboxedalgorithm}{%
	%
	enhanced,
	sharp corners,
	frame hidden,
	%
	center,
	width=\linewidth,
	top=0.5em,
	bottom=0.5em,
	left=0.5em,
	right=0.5em,
	%
	before skip=0.5em,
	after skip=0.5em,
	%
	attach title to upper,
	coltitle=black,
	%
	breakable
}


% List-like algorithm setup.
\newlist{enumalgorithm}{enumerate}{2}
\setlist[enumalgorithm,1]{nosep, topsep=0.25em, itemsep=2.5pt, label=(\arabic*), leftmargin=*}
\setlist[enumalgorithm,2]{topsep=2.5pt, parsep=0pt, itemsep=2.5pt, label=(\alph*), ref=(\arabic{enumalgorithmi}\alph*)}


% Alternate description environment.
\newlist{descriptor}{description}{1}
\setlist[descriptor]{labelindent=0.25in,labelwidth=0.25in,leftmargin=!,align=left,itemsep=0pt,topsep=0.5em,rightmargin=\leftmargin,font={\bf\textit}}

\newlist{widedescriptor}{description}{1}
\setlist[widedescriptor]{topsep=0pt,labelindent=0.25in,labelwidth=0.25in,leftmargin=!,align=left,itemsep=0pt,topsep=0.5em,font={\bf\textit}}


% Inset algorithm environment that wraps around boxedalgorithm and enumalgorithm.
\makeatletter

\define@key{insetalgorithm@keys}{name}{\def\insetalgorithm@name{#1}}%
\define@key{insetalgorithm@keys}{label}{\def\insetalgorithm@label{#1}}%
\define@key{insetalgorithm@keys}{description}{\def\insetalgorithm@description{#1}}%

\newenvironment{insetalgorithm}[1][]
	{%
		\setkeys{insetalgorithm@keys}{name={},label=algorithm,description={},#1}
		\begin{boxedalgorithm}{
			\ifthenelse{\equal{\insetalgorithm@name}{}}{}{%
				\textit{(\insetalgorithm@name)}%
			}.%
		}{\insetalgorithm@label}{%
			\ifthenelse{\equal{\insetalgorithm@description}{}}{\\[-0.75em]%
			}{%
                \\[0.5em]\insetalgorithm@description%
			}
		}
			\begin{enumalgorithm}%
	}
	{%
			\end{enumalgorithm}
		\end{boxedalgorithm}%
	}

\makeatother



%%%%%%%%%%%%%%%%%%%%%%%%%%%%%%%%%%%%%%%%%%%%%%%%%%%%%%%%%%%%%%%%%%%%%%%%%%%%%%%%
%%%% BIBLIOGRAPHY                                                           %%%%
%%%%%%%%%%%%%%%%%%%%%%%%%%%%%%%%%%%%%%%%%%%%%%%%%%%%%%%%%%%%%%%%%%%%%%%%%%%%%%%%
\usepackage{hyperref}
\usepackage[
	backend=biber,
	style=numeric-comp,
	sorting=noneyear,
	sortcites,
	maxnames=3,
	date=year
]{biblatex}

% Italicized titles
\DeclareFieldFormat[article,inproceedings,book,online,incollection]{title}{\emph{#1}}

% Remove extraneous fields.
\AtEveryBibitem{\clearfield{pages}}
\AtEveryBibitem{\clearfield{volume}}
\AtEveryBibitem{\clearfield{publisher}}
\AtEveryBibitem{\clearfield{editor}}
\AtEveryBibitem{\clearfield{series}}
\AtEveryBibitem{\clearfield{number}}
\AtEveryBibitem{\clearfield{isbn}}

\AtEveryBibitem{\clearlist{series}}
\AtEveryBibitem{\clearlist{publisher}}
\AtEveryBibitem{\clearlist{editor}}

% Re-format useful fields.
\DeclareFieldFormat*{url}{}
\DeclareFieldFormat[software]{url}{\mkbibacro{URL}\addcolon\space\url{#1}}

\DeclareFieldFormat*{doi}{}
\DeclareFieldFormat[online,software]{doi}{\mkbibacro{DOI}\addcolon\space\url{#1}}


% Smaller (list) label sep, itemsep, font size, multiple citations.
\setlength{\biblabelsep}{\labelsep}
\renewcommand*{\bibitemsep}{0pt}
\renewcommand*{\bibfont}{\small}
\renewcommand*{\multicitedelim}{\addcomma\nobreakspace}

% Sort references by year within the citation.
\DeclareSortingTemplate{noneyear}{
	\sort{\citeorder}
	\sort{\field{year}}
}

\addbibresource{main.bib}



%%%%%%%%%%%%%%%%%%%%%%%%%%%%%%%%%%%%%%%%%%%%%%%%%%%%%%%%%%%%%%%%%%%%%%%%%%%%%%%%
%%%% SPECIAL MATHEMATICS TYPESETTING                                        %%%%
%%%%%%%%%%%%%%%%%%%%%%%%%%%%%%%%%%%%%%%%%%%%%%%%%%%%%%%%%%%%%%%%%%%%%%%%%%%%%%%%
\usepackage{amsmath, amssymb, nicefrac, mathtools, bm}

\renewcommand{\qed}{\hfill$\blacksquare$}

% Allow aligned equations to break across pages.
\allowdisplaybreaks

\let\epsilon\varepsilon

\providecommand{\N}{\mathds{N}}					% naturals
\providecommand{\F}{\mathds{F}}					% arbitrary field
\providecommand{\R}{\mathds{R}}					% reals
\providecommand{\C}{\mathds{C}}					% complex
\providecommand{\T}{\mathds{T}}					% torus?
\providecommand{\Z}{\mathds{Z}}					% integers
\providecommand{\Q}{\mathds{Q}}					% rationals
\providecommand{\RP}{\mathds{R}\mathrm{P}}		% projective plane

\let\P\undefined
\let\H\undefined

\DeclareMathOperator{\P}{\mathbf{P}}				% probability measure
\DeclareMathOperator{\E}{\mathbf{E}}				% expectation
\DeclareMathOperator{\1}{\mathbf{1}}				% indicator
\DeclareMathOperator{\Uniform}{\mathbf{Uniform}}	% uniform distribution

\let\im\undefined
\DeclareMathOperator{\im}{\mathrm{im}}			% algebraic image
\DeclareMathOperator{\Hom}{\mathrm{Hom}}		% Hom functor
\DeclareMathOperator{\Ext}{\mathrm{Ext}}		% Ext functor
\DeclareMathOperator{\Tor}{\mathrm{Tor}}		% Tor functor
\DeclareMathOperator{\rank}{\mathrm{rank}}		% algebraic rank
\DeclareMathOperator{\Range}{\mathrm{Range}}	% range

\DeclareMathOperator{\Length}{\mathrm{Length}}		% length
\DeclareMathOperator{\Area}{\mathrm{Area}}			% area
\DeclareMathOperator{\Volume}{\mathrm{Vol}}			% volume
\DeclareMathOperator{\Perimeter}{\mathrm{Perim}}	% perimeter
\DeclareMathOperator{\Trees}{\mathrm{Trees}}

\DeclareMathOperator{\Volumizer}{\mathcal V}
\DeclareMathOperator{\Perimetrizer}{\mathcal P}

\newcommand{\inverse}[1]{{#1}^{-1}}		% inverse

\DeclareMathOperator{\Link}{\mathrm{Link}}		% linking number
\newcommand{\Mod}[1]{\ (\mathrm{mod}\ #1)}		% better modulus

\newcommand{\Reduced}{\widetilde{H}}						% reduced homology
\newcommand{\ReducedPersistent}{\widetilde{\mathit{PH}}}	% reduced persistent homology

\newcommand{\Normalizer}{\mathcal{Z}}		% normalization constant

\newcommand{\dualforest}[1]{{#1}^\ast}		% dual forest
\newcommand{\dual}[1]{{#1}^\bullet}			% dual subcomplex

\DeclareMathOperator{\Alexander}{\mathcal{A}}					% Alexander duality
\DeclareRobustCommand{\rAlexander}{\reflectbox{$\mathcal A$}}	% backwards cal A
\DeclareMathOperator{\Dualize}{\mathcal{D}}						% dualization map

\DeclareMathOperator{\Int}{\mathrm{Int}}	% topological interior
\DeclareMathOperator{\Bd}{\mathrm{Bd}}		% topological boundary
\DeclareMathOperator{\Cl}{\mathrm{Cl}}		% topological closure

\let\into\hookrightarrow		% injective function
\let\onto\twoheadrightarrow		% surjective function
\let\lto\longrightarrow			% longer arrow (it's prettier)
\let\lmapsto\longmapsto         % longer mapsto arrow

\DeclareMathOperator{\CPP}{\bm{\mathsf{CPP}}}		% coupled plaquette percolation
\DeclareMathOperator{\UST}{\bm{\mathsf{UST}}}       % uniform spanning tree distribution
\DeclareMathOperator{\USF}{\bm{\mathsf{USF}}}       % uniform spanning forest distribution
\DeclareMathOperator{\PLGT}{\bm{\mathsf{PLGT}}}     % Potts lattice gauge theory
\DeclareMathOperator{\PRCM}{\bm{\mathsf{PRCM}}}     % Plaquette random cluster model
\DeclareMathOperator{\PK}{\bm{\mathsf{PK}}}         % Coupling

\DeclareMathOperator{\H}{\mathbf{H}}        % Hamiltonian

\newcommand{\skeleton}[2]{{#1}^{(#2)}}		% skeleton

\let\complement\overline	% complement



%%%%%%%%%%%%%%%%%%%%%%%%%%%%%%%%%%%%%%%%%%%%%%%%%%%%%%%%%%%%%%%%%%%%%%%%%%%%%%%%
%%%% SPECIAL NONMATHEMATICS TYPESETTING                                     %%%%
%%%%%%%%%%%%%%%%%%%%%%%%%%%%%%%%%%%%%%%%%%%%%%%%%%%%%%%%%%%%%%%%%%%%%%%%%%%%%%%%
\usepackage{xspace}

\let\th\undefined
\let\st\undefined
\let\nd\undefined
\let\rd\undefined

\newcommand{\th}{\textsuperscript{th}\xspace}
\newcommand{\st}{\textsuperscript{st}\xspace}
\newcommand{\nd}{\textsuperscript{nd}\xspace}
\newcommand{\rd}{\textsuperscript{rd}\xspace}

\newcommand{\qand}{\quad \text{and} \quad}
\newcommand{\qbut}{\quad \text{but} \quad}
\newcommand{\qor}{\quad \text{or} \quad}

\newcommand{\qqand}{\qquad \text{and} \qquad}
\newcommand{\qqbut}{\qquad \text{but} \qquad}
\newcommand{\qqor}{\qquad \text{or} \qquad}

\newcommand{\ATEAMS}{\textsf{ATEAMS}\xspace}
\newcommand{\ATEAMSpp}{\textsf{ATEAMS++}\xspace}



%%%%%%%%%%%%%%%%%%%%%%%%%%%%%%%%%%%%%%%%%%%%%%%%%%%%%%%%%%%%%%%%%%%%%%%%%%%%%%%%
%%%% MISC                                                                   %%%%
%%%%%%%%%%%%%%%%%%%%%%%%%%%%%%%%%%%%%%%%%%%%%%%%%%%%%%%%%%%%%%%%%%%%%%%%%%%%%%%%

% https://tex.stackexchange.com/a/59769
% Patch to reduce space before/after proofs
\makeatletter
\renewenvironment{proof}[1][\proofname]{\par
	\vspace{-\topsep}% remove the space after the theorem
	\pushQED{\qed}%
	\normalfont
	\topsep0.5em \partopsep0pt %
	\trivlist
	\item[\hskip\labelsep
	    \itshape
	#1\@addpunct{.}]\ignorespaces
}{%
	\popQED\endtrivlist\@endpefalse
	\addvspace{0.25em plus 2pt} % some space after
}
\makeatother

% Reduce spacing around displaymath that DOESN'T get reset every time.
\makeatletter
\g@addto@macro \normalsize {%
	\setlength\abovedisplayskip{5pt plus 2pt minus 2pt}%
	\setlength\belowdisplayskip{5pt plus 2pt minus 2pt}%
}
\makeatother

